\section{创新假设、风险与下一轮决定}\label{sec:decisions}
候选方法创新是：（1）BFZI 匝号仿射结构同时用于
积分核及神经场源编码；（2）具有海水电流连续约束的非局部场响应学习，
用同一响应计算端口、空间场、损耗；
（3）随体固体温度与世界海水尾迹的双坐标表示；
（4）同源功率一致的热降阶 ODE 与谐振电路反馈。
POD、物理先验与 NN、参数化热 ROM、LLC 电热代理、
神经算子与双向磁热耦合已有先例
\cite{wang2026,dam2026,zhu2026,xie2026,yuan2024,chen2025,dini2023,barba2026}。
仅将这些名词组合不足以证明创新，必须以明确计算复杂度、
误差、跨几何消融和实验结果建立独立贡献。

\textbf{不可声称}任意材料/任意高速运动无需验证、
无固定世界网格等同完全免求积、
端口准确等同空间热源准确、
单频被动性保证全频因果性、
已有论文的速度直接可迁移为本模型速度。

代码实现前应收敛下列五个具体决定：
完整端口和导体截面约定；近场集肤/邻近的基准方法；
海水和导电盒壁材料界面表示；低维电磁状态/投影设计；
几何条件化热基如何保留耗散与跨界面条件。
新理论保留与实验事实的区别，后续以本目录为主线，
不再以旧超椭圆理论并行指导 BFZI 代码。
